% !TEX root = ../main.tex

\begin{acknowledgements}
We would like to express our sincere gratitude to our mentor, Prof. Yutong Ban, for his guidance throughout this project. His advice on research direction, robotic-system integration, experimental design, and technical presentation was essential to the development of the dual-arm manipulation platform and the completion of this thesis. His careful feedback encouraged us to examine the system objectively, distinguish demonstrated results from unsupported claims, and continuously improve both the engineering work and its documentation.

We are also deeply grateful to our instructor, Prof. Peisen Huang, for his guidance during the design-review process. His instruction on translating customer requirements into measurable engineering specifications, evaluating design alternatives, validating prototype performance, and communicating engineering decisions provided the methodological foundation for this project.

We sincerely thank SJTU SIRIUS Lab for providing the experimental platform, equipment, computing resources, and technical environment required for this work. We are especially grateful to Han Fang for his patient assistance and valuable guidance throughout the development, integration, and experimental validation of the system. His support helped us address practical difficulties encountered during the project and contributed significantly to the successful completion of the experiments.

We thank the faculty and staff of Global College, Shanghai Jiao Tong University, for providing the academic environment and organizational support necessary for this work. We also appreciate the Student Innovation Center for providing fabrication facilities and assistance with the custom grippers, fixtures, and other prototype components.

We further acknowledge the developers and research communities behind OpenPI, LeRobot, GELLO, ROS~2, PyTorch, and the related open-source tools used in this project. Their publicly available software, documentation, and research made it possible to construct, train, and deploy the integrated $\pi_{0.5}$ system within the project schedule.

Finally, we would like to thank our families and friends for their patience, encouragement, and support throughout the project. Their understanding sustained us during the intensive periods of system integration, data collection, model training, physical validation, and thesis preparation.
\end{acknowledgements}
